% pre-v1.tex
%
% Purpose:
%   Beamer presentation of the theory part (sections 1--5 and appendix A)
%   of the rewritten draft on guided modes in line-defect photonic-crystal
%   waveguides: model, Cauchy-data center reduction, BIE / Rayleigh-port
%   discretization, and the resolvent-based spectral distance bound.
%
% Main workflow:
%   model and spectrum -> center reduction via half-guide Cauchy-data spaces
%   -> Rayleigh ports, layer potentials, one-cell scattering, ordered QZ
%   -> discrete NEP M(k) and candidate scan -> raw field, interface jumps,
%   trace estimate, spectral distance theorem, computable / certified bounds.
%
% Based on:
%   draft/tmp/theory-rewrite/{s-introduction,s-modelproblem,s-reduction,
%   s-discretization,s-error,a-cellrep}.tex, and Fliss (2013),
%   Joly--Li--Fliss (2006) for the continuous half-guide theory.
%
% Main differences:
%   Unlike pre1--pre3, this deck follows the rewritten draft notation
%   (M(k), E_L/E_R, U_L/U_R, omega for the single-layer density, L_0 for the
%   center width) and uses the Cauchy-data formulation as the main line.
%   Numerical results are intentionally not included.
%
% Numerical goal:
%   None; this deck states the theory needed before presenting numerical
%   spectral distance certificates.

\documentclass[aspectratio=169]{beamer}

\usetheme{CambridgeUS}
\usecolortheme{dove}
\usefonttheme{professionalfonts}

\usepackage{amsmath}
\usepackage{amssymb}
\usepackage{mathrsfs}
\usepackage{mathtools}
\usepackage{stmaryrd}
\usepackage{bm}
\usepackage{booktabs}
\usepackage{tikz}
\usetikzlibrary{arrows.meta,calc,decorations.pathreplacing,positioning}

\renewcommand{\Im}{\mathrm{Im}}
\renewcommand{\Re}{\mathrm{Re}}
\newcommand{\ii}{\mathrm{i}}
\newcommand{\ee}{\mathrm{e}}
\newcommand{\dd}{\mathrm{d}}
\newcommand{\CC}{\mathbb{C}}
\newcommand{\RR}{\mathbb{R}}
\newcommand{\ZZ}{\mathbb{Z}}
\newcommand{\jump}[1]{\llbracket #1\rrbracket}
\newcommand{\QP}{\mathrm{QP}}
\newcommand{\rin}{\rho_{\min}}
\newcommand{\rax}{\rho_{\max}}
\newcommand{\eps}{\varepsilon_{\partial,\ell_*}}
\newcommand{\dist}{\operatorname{dist}}

\newtheorem{proposition}[theorem]{Proposition}
\newtheorem{assumption}[theorem]{Assumption}
\newtheorem{remark}[theorem]{Remark}

\setbeamertemplate{navigation symbols}{}

\title[Guided modes: BIE + spectral distance]{Guided Modes in Line-Defect Periodic Waveguides:\\
  Cauchy-Data Reduction, BIE Discretization,\\and Spectral Distance Bounds}
\subtitle{Theory part (pre-v1)}
\author{}
\date{}

\AtBeginSection[]{
  \begin{frame}{Outline}
    \tableofcontents[currentsection]
  \end{frame}
}

\begin{document}

\begin{frame}
  \titlepage
\end{frame}

\begin{frame}{Outline}
  \tableofcontents
\end{frame}

% ======================================================================
\section{Problem and Goals}

\begin{frame}{Motivation}
  \small
  A line defect in a 2D photonic crystal breaks periodicity across the defect
  but keeps it along the defect. Inside a band gap it may support
  \textbf{guided modes}: waves propagating along the defect and exponentially
  localized across it \cite{JoannopoulosJohnsonWinnMeade2008}.

  \medskip
  For fixed Bloch number $\beta$, guided modes are eigenvalues of a
  self-adjoint operator $A(\beta)$ on an unbounded strip, lying in a gap of
  the essential spectrum \cite{Fliss2013}.

  \begin{block}{Two difficulties}
    \begin{itemize}
      \item \textbf{Unbounded transverse direction.} Supercells introduce a model
        error. Exact half-guide conditions
        \cite{JolyLiFliss2006,Fliss2013} reduce the problem to a bounded region,
        at the price of a nonlinear eigenproblem in $k$.
      \item \textbf{Reliability.} BIE methods are very efficient for
        piecewise-constant media
        \cite{Linton1998,BarnettGreengard2010,HaoBarnettMartinssonYoung2014}, but a
        nearly singular matrix only says a finite system is nearly solvable
        \cite{GuettelTisseur2017}. It does not bound the distance to
        $\sigma(A(\beta))$.
    \end{itemize}
  \end{block}
\end{frame}

\begin{frame}{This work}
  \footnotesize
  \begin{enumerate}
    \item \textbf{Center reduction via Cauchy data.} Guided modes $\Longleftrightarrow$
      nonzero center solutions whose Cauchy data lie in the half-guide Cauchy-data
      spaces $\mathcal G_\pm(k)$. No Dirichlet well-posedness is needed. DtN is
      the special case where $\mathcal G_\pm$ is a graph.
    \item \textbf{BIE + Rayleigh-port discretization.} Layer potentials on
      inclusions, Rayleigh amplitudes on ports, one-cell scattering, and
      ordered QZ give the decaying subspace. The result is a discrete NEP
      $\mathsf M(k)$ in Cauchy-matching form that produces candidates $\widehat k$.
    \item \textbf{Spectral distance theorem.} A piecewise field $\widehat u$ built
      from the candidate solves Helmholtz exactly in each piece. Its interface jumps
      control $\dist(\widehat k^2,\sigma(A(\beta)))$, with no prior gap verification.
    \item \textbf{Computable bounds.} Material-interface and port integrals, a
      propagation-based tail for infinitely many cells, and one-sided
      (certified) inputs.
  \end{enumerate}

  \begin{alertblock}{Not claimed}
    The bound only says that an interval around $\widehat k^2$ meets
    $\sigma(A(\beta))$. Calling the spectral point a guided-mode eigenvalue
    needs a separate proof that the interval lies in a gap. No uniqueness,
    no indexing, and no a priori convergence of the discretization are claimed.
  \end{alertblock}
\end{frame}

\begin{frame}{Geometry}
  \centering
  \begin{tikzpicture}[x=0.95cm,y=0.95cm,
      every node/.style={font=\scriptsize},
      incl/.style={draw=blue!60!black, fill=blue!18, line width=0.5pt},
      port/.style={draw=red!70!black, line width=1.1pt},
      art/.style={draw=black!55, dashed, line width=0.5pt}]
    \fill[gray!6] (-6.4,-1) rectangle (6.4,1);
    \fill[orange!12] (-1,-1) rectangle (1,1);
    \draw[thick] (-6.4,1) -- (6.4,1);
    \draw[thick] (-6.4,-1) -- (6.4,-1);
    \foreach \x in {-5.8,-4.2,-2.6,2.6,4.2,5.8} \draw[art] (\x,-1) -- (\x,1);
    \draw[port] (-1,-1) -- (-1,1);
    \draw[port] (1,-1) -- (1,1);
    \foreach \c in {1.75,3.35,4.95} \draw[incl] (\c,0.05) circle (0.42);
    \foreach \c in {-1.85,-3.45,-5.05} \draw[incl] (\c,-0.05) circle (0.42);
    \draw[draw=green!45!black, fill=green!20, line width=0.5pt] (0,0) ellipse (0.55 and 0.3);
    \node at (0,0) {$\Omega_0$};
    \node at (1.75,0.05) {$\Omega_+$};
    \node at (-1.85,-0.05) {$\Omega_-$};
    \node at (0,0.7) {$\mathcal C_0$};
    \node at (1.8,0.75) {$\mathcal C_{0,+}$};
    \node at (3.4,0.75) {$\mathcal C_{1,+}$};
    \node at (5.0,0.75) {$\mathcal C_{2,+}$};
    \node at (-1.8,0.75) {$\mathcal C_{0,-}$};
    \node at (-3.4,0.75) {$\mathcal C_{1,-}$};
    \node at (-5.0,0.75) {$\mathcal C_{2,-}$};
    \node at (6.7,0) {$\cdots$};
    \node at (-6.7,0) {$\cdots$};
    \node[below] at (-1,-1) {$\Gamma_-$};
    \node[below] at (1,-1) {$\Gamma_+$};
    \draw[<->, black!60] (-1,-1.55) -- (1,-1.55) node[midway, below] {$L_0$};
    \draw[<->, black!60] (1,-1.55) -- (2.6,-1.55) node[midway, below] {$L$};
    \node[below] at (2.6,-1) {$\Gamma_{0,+}$};
    \node[below] at (4.2,-1) {$\Gamma_{1,+}$};
    \node[below] at (-2.6,-1) {$\Gamma_{0,-}$};
    \node[below] at (-4.2,-1) {$\Gamma_{1,-}$};
    \draw[-{Latex[length=1.6mm]}, red!70!black] (1,-0.75) -- (1.4,-0.75)
      node[right, inner sep=1pt] {$\nu_+$};
    \draw[-{Latex[length=1.6mm]}, red!70!black] (-1,-0.75) -- (-1.4,-0.75)
      node[left, inner sep=1pt] {$\nu_-$};
    \draw[decorate, decoration={brace, amplitude=4pt}] (1,1.08) -- (6.4,1.08)
      node[midway, above=3pt] {half-guide $B_+$ (periodic $\rho_{\mathrm p}$, period $L$)};
    \draw[decorate, decoration={brace, amplitude=4pt}] (-6.4,1.08) -- (-1,1.08)
      node[midway, above=3pt] {half-guide $B_-$};
    \draw[decorate, decoration={brace, amplitude=4pt}] (-1,1.08) -- (1,1.08)
      node[midway, above=3pt] {defect};
    \draw[<->, black!60] (-7.0,-1) -- (-7.0,1) node[midway, left] {$d$};
  \end{tikzpicture}

  \vspace{0.4em}
  \raggedright\footnotesize
  Strip $B=\RR\times I_d$, $I_d=(-d/2,d/2)$, $\beta$-quasi-periodic in $y$;
  ports at $x_\pm=\pm L_0/2$.
  Background $\rho_{\mathrm e}$; one $C^2$ simply connected inclusion per cell
  ($\rho_{\mathrm c}$; center: $\rho_{0,\mathrm c}$ or empty). Each inclusion's closure
  lies inside its cell. $\Omega_\pm$ are translates of the same inclusion,
  possibly at different relative positions.
\end{frame}

\begin{frame}{Guided modes as eigenvalues of $A(\beta)$}
  \small
  Guided mode: a nonzero solution of $-\Delta w=k^2\rho w$ in $\RR^2$ with
  $w=v\,\ee^{\ii\beta y}$, $v$ $d$-periodic in $y$, and $v|_B\in H^1(B)$
  (transverse localization).

  \medskip
  With $\mathcal H:=L^2(B,\rho\,\dd x\,\dd y)$ and $0<\rin\le\rho\le\rax$, define
  \cite{Fliss2013}
  \[
    A(\beta)u:=-\rho^{-1}\Delta u,\qquad
    D(A(\beta)):=\bigl\{u\in H^1_\beta(B):\ \Delta u\in L^2(B),\
      \partial_yu|_{y=d/2}=\ee^{\ii\beta d}\partial_yu|_{y=-d/2}\bigr\}.
  \]
  \begin{itemize}
    \item $\Delta u$ is distributional, so $\Delta u\in L^2$ already enforces
      continuity of $u$ and $\partial_\nu u$ across material interfaces.
    \item $A(\beta)$ is nonnegative self-adjoint, and
      $(A(\beta)u,u)_{\mathcal H}=\lVert\nabla u\rVert_{L^2(B)}^2$.
  \end{itemize}

  \begin{block}{Eigenproblem}
    Find $k^2>0$ and $u\in D(A(\beta))\setminus\{0\}$ such that $A(\beta)u=k^2u$.
  \end{block}
\end{frame}

\begin{frame}{Spectrum}
  \small
  \begin{proposition}[Essential spectrum {\cite[Prop.~3.1]{Fliss2013}}]
    \[
      \sigma_{\mathrm{ess}}(A(\beta))=\sigma(A_{\mathrm p}(\beta))
      =\bigcup_{\theta\in[-\pi/L,\pi/L)}\sigma(A_{\mathrm p}(\beta,\theta))
      =[0,\infty)\setminus\bigcup_{n=1}^{N(\beta)}(a_n(\beta),b_n(\beta)).
    \]
  \end{proposition}
  \vspace{-0.3em}
  {\footnotesize The defect is a compact resolvent perturbation; the rest is
  Floquet--Bloch. Fliss only needs $\rho\in L^\infty$ with positive bounds, so
  piecewise-constant $\rho$ is included.}

  \begin{assumption}[Gap]
    Fix $\beta$ and a gap $G_\beta=(g_-,g_+)$, $0<g_-<g_+<\infty$, with
    $G_\beta\cap\sigma_{\mathrm{ess}}(A(\beta))=\varnothing$.
  \end{assumption}

  \begin{proposition}[Discrete spectrum in the gap {\cite[Prop.~3.3, Thm.~3.5]{Fliss2013}}]
    $\sigma(A(\beta))\cap G_\beta$ consists of isolated eigenvalues of finite
    multiplicity, accumulating at most at $g_\pm$. Their eigenfunctions decay
    exponentially as $|x|\to\infty$.
  \end{proposition}
  {\footnotesize Non-emptiness is not guaranteed. Candidates are searched in
  $G_\beta$; the distance theorem holds for all of $\sigma(A(\beta))$.}
\end{frame}

% ======================================================================
\section{Center Reduction via Cauchy Data}

\begin{frame}{Cauchy data spaces of the half-guides}
  \small
  Cells: $\mathcal C_0=(x_-,x_+)\times I_d$ and $\mathcal C_{j,\pm}=\mathcal C_\pm\pm jLe_x$.
  Ports $\Gamma_\pm=\{x_\pm\}\times I_d$ with $\nu_\pm=\pm e_x$ pointing \emph{into}
  $B_\pm$. All normal derivatives on $\Gamma_\pm$ are taken along $\nu_\pm$.

  \medskip
  Cauchy map
  $\mathscr T_\pm u:=\bigl(\gamma_\pm u,\ \partial_{\nu_\pm}u|_{\Gamma_\pm}\bigr)\in
  H^{1/2}_\beta(\Gamma_\pm)\times H^{-1/2}_\beta(\Gamma_\pm)$. Local solution spaces
  (distributional equation, \emph{no} condition on $\Gamma_\pm$):
  \begin{align*}
    \mathcal U_0(k)&:=\{u\in H^1_\beta(\mathcal C_0):\ -\Delta u=k^2\rho u\ \text{in }\mathcal C_0,\
      \partial_yu\ \beta\text{-q.p.}\},\\
    \mathcal V_\pm(k)&:=\{v\in H^1_\beta(B_\pm):\ -\Delta v=k^2\rho_{\mathrm p}v\ \text{in }B_\pm,\
      \partial_yv\ \beta\text{-q.p.}\}.
  \end{align*}

  \begin{block}{Half-guide Cauchy data space}
    \[
      \mathcal G_\pm(k):=\mathscr T_\pm\mathcal V_\pm(k)
      \subset H^{1/2}_\beta(\Gamma_\pm)\times H^{-1/2}_\beta(\Gamma_\pm).
    \]
    It collects the Cauchy data of \emph{all} finite-energy solutions in $B_\pm$.
  \end{block}
\end{frame}

\begin{frame}{Unique continuation and the center reduction}
  \small
  \begin{lemma}[Unique continuation]
    If $v\in\mathcal V_\pm(k)$ and $\mathscr T_\pm v=0$, then $v=0$.
  \end{lemma}
  {\footnotesize Proof idea: near $\Gamma_\pm$ the coefficient is $\rho_{\mathrm e}$.
  Extend by zero, use analyticity in the connected background, then the transmission
  conditions and the Green representation inside the inclusions.}

  \begin{proposition}[Center reduction]
    Fix $\beta$ and $k>0$, and let
    \[
      \mathcal U_0^{\mathcal G}(k):=\bigl\{u_0\in\mathcal U_0(k):\
        \mathscr T_-u_0\in\mathcal G_-(k),\ \mathscr T_+u_0\in\mathcal G_+(k)\bigr\}.
    \]
    The restriction $u\mapsto u|_{\mathcal C_0}$ is a linear bijection
    $\ker(A(\beta)-k^2)\to\mathcal U_0^{\mathcal G}(k)$. In particular, $k^2$ is an
    eigenvalue if and only if $\mathcal U_0^{\mathcal G}(k)\ne\{0\}$, with equal
    dimensions.
  \end{proposition}
\end{frame}

\begin{frame}{Center reduction: proof}
  \small
  \textbf{Proof.}
  \begin{itemize}
    \item \emph{Well defined:} $u\in H^1$ with $\Delta u\in L^2$, so the value trace and
      the $\nu_\pm$-normal trace agree on both sides of $\Gamma_\pm$.
    \item \emph{Injective:} $u|_{\mathcal C_0}=0$ gives $\mathscr T_\pm(u|_{B_\pm})=0$,
      so $u|_{B_\pm}=0$ by the lemma.
    \item \emph{Surjective:} pick $v_\pm\in\mathcal V_\pm$ with
      $\mathscr T_\pm v_\pm=\mathscr T_\pm u_0$ and glue $v_-,u_0,v_+$. Matching
      value traces give $H^1_\beta(B)$. Matching normal traces mean $\Delta u$ has no
      singular part on $\Gamma_\pm$.
  \end{itemize}

  \begin{alertblock}{}
    This holds for \emph{every} $k>0$. It needs neither the gap assumption nor
    unique solvability of a half-guide Dirichlet problem.
  \end{alertblock}
\end{frame}

\begin{frame}{Relation to DtN and to the propagation operator}
  \small
  \textbf{DtN as a special case.} Suppose every $\varphi\in H^{1/2}_\beta(\Gamma_\pm)$ has a
  unique $v\in\mathcal V_\pm(k)$ with $\gamma_\pm v=\varphi$. Then
  \[
    \mathcal G_\pm(k)=\operatorname{graph}\Lambda_\pm(k),\qquad
    \Lambda_\pm(k)\varphi:=\partial_{\nu_\pm}v|_{\Gamma_\pm},
  \]
  and $\mathscr T_\pm u_0\in\mathcal G_\pm$ becomes the transparent condition
  $\partial_{\nu_\pm}u_0=\Lambda_\pm\gamma_\pm u_0$. This is Fliss' reduction
  \cite[Thm.~4.5]{Fliss2013}; her DtN, taken with the outward normal of $B_\pm$,
  equals $-\Lambda_\pm$.

  \begin{itemize}
    \item In the gap, the half-guide Dirichlet problem is well posed except on a
      countable set \cite[Thm.~4.1]{Fliss2013}: the eigenvalues of the half-guide
      operator with Dirichlet data on $\Gamma_\pm$.
    \item At those values $\mathcal G_\pm$ is not a graph and DtN fails, but the
      Cauchy-data reduction still holds. (Robin-to-Robin is another workaround
      \cite{FlissKlindworthSchmidt2015}.)
  \end{itemize}

  \textbf{Structure of $\mathcal V_\pm(k)$ in the gap.} There are no $x$-propagating
  Bloch waves. Finite-energy solutions decay cell by cell through a propagation
  operator $\mathcal R_\pm$ with spectral radius $<1$ \cite{JolyLiFliss2006}
  (Jordan structure: \cite{HohageSoussi2013}).\\
  $\Rightarrow$ discrete counterpart: the cell propagation matrix $\mathsf P_\pm$.
\end{frame}

% ======================================================================
\section{Discretization and Candidates}

\begin{frame}{Continuous objects and their discrete counterparts}
  \small
  \centering
  \renewcommand{\arraystretch}{1.25}
  \begin{tabular}{lll}
    \toprule
    continuous & discrete & built from\\
    \midrule
    $\mathcal V_\pm(k)$, operator $\mathcal R_\pm$ &
      $\mathsf Z_\pm$, $\mathsf P_\pm$ & one-cell scattering + ordered QZ\\
    $\mathcal G_\pm(k)=\mathscr T_\pm\mathcal V_\pm(k)$ &
      $\operatorname{range}\mathsf V_\pm$, $\mathsf V_\pm=\mathsf C_\pm\mathsf Z_\pm$ &
      amplitude $\to$ Cauchy data\\
    $\mathcal U_0(k)$ &
      $\xi=(\eta;a^-_0;b^+_0)$ & layer potentials + $u_{\mathrm{hom}}$\\
    transmission on $\partial\Omega_0$ & $J(k)\xi=0$ & Nystr\"om jump operator\\
    $\mathscr T_\pm u_0\in\mathcal G_\pm(k)$ &
      $\Pi_\pm\xi=\mathsf V_\pm c_\pm$ & port Rayleigh coefficients\\
    $\ker(A(\beta)-k^2)\ne\{0\}$ & $\mathsf M(k)$ singular & candidate scan\\
    \bottomrule
  \end{tabular}

  \vspace{0.8em}
  \raggedright
  Fix $\beta$ and take $k\in(\sqrt{g_-},\sqrt{g_+})$. The candidates $\widehat k$ are only
  \emph{inputs} to the spectral distance estimate; this section makes no accuracy
  claim.
\end{frame}

\begin{frame}{Port Fourier basis and Rayleigh waves}
  \footnotesize
  $\beta_m:=\beta+2\pi m/d$ and $\psi_m(y):=d^{-1/2}\ee^{\ii\beta_my}$, truncated to $|m|\le M$
  ($K=2M+1$). With $\kappa_{\mathrm e}=k\sqrt{\rho_{\mathrm e}}$:
  \[
    \gamma_m:=\sqrt{\kappa_{\mathrm e}^2-\beta_m^2},\quad\Im\gamma_m\ge0\
    (\gamma_m\ge0\text{ if real}),\qquad
    \boldsymbol\gamma:=\operatorname{diag}(\gamma_m).
  \]
  \textbf{Assumption (no Wood anomaly):} $\gamma_m(k)\neq0$ for all $m$. (At
  $\gamma_m=0$, $G_\QP$ blows up \cite{BarnettGreengard2010}.)

  \medskip

  In a background strip with anchor $x_p$:
  \[
    r_m^{x_p}:=\ee^{\ii\gamma_m(x-x_p)}\psi_m,\qquad
    l_m^{x_p}:=\ee^{-\ii\gamma_m(x-x_p)}\psi_m,\qquad
    u=\sum_m(a_mr_m^{x_p}+b_ml_m^{x_p}).
  \]
  \begin{itemize}
    \item $(a,b)$ are the right-/left-going amplitudes. Moving the anchor by $\ell$
      multiplies them by $\mathsf E(\ell):=\operatorname{diag}(\ee^{\ii\gamma_m\ell})$.
    \item Amplitudes $\to$ Cauchy-data coefficients for port normal $\pm e_x$:
      \[
        \mathsf C_\pm:=\begin{bmatrix}I&I\\ \pm\ii\boldsymbol\gamma&\mp\ii\boldsymbol\gamma\end{bmatrix}
        \in\CC^{2K\times2K},\qquad\text{invertible under no-Wood.}
      \]
    \item Natural weights: $(1+\beta_m^2)^{\pm1/4}$ for value ($H^{1/2}$) and normal
      derivative ($H^{-1/2}$) coefficients.
  \end{itemize}
\end{frame}

\begin{frame}{Layer potentials and the transmission BIE}
  \footnotesize
  $G^{(\kappa)}=\tfrac{\ii}{4}H_0^{(1)}(\kappa|\mathbf z-\mathbf z'|)$, and for $x\neq x'$
  \cite{Linton1998}
  \[
    G^{(\kappa)}_\QP(\mathbf z,\mathbf z')=\frac{\ii}{2d}\sum_{m\in\ZZ}\frac1{\gamma_m}
    \ee^{\ii\gamma_m|x-x'|+\ii\beta_m(y-y')}.
  \]
  Cell representation, with $\nu$ pointing from inside to outside and
  $\kappa_{\mathrm c}=k\sqrt{\rho_{\mathrm c}}$:
  \[
    u=\begin{cases}
      \mathcal S^{(\kappa_{\mathrm e})}_\QP[\omega]+\mathcal D^{(\kappa_{\mathrm e})}_\QP[\tau]+u_{\mathrm b},
        & \text{outside }\Omega,\\[0.2em]
      \mathcal S^{(\kappa_{\mathrm c})}[\omega]+\mathcal D^{(\kappa_{\mathrm c})}[\tau], & \text{inside }\Omega,
    \end{cases}
    \qquad u_{\mathrm b}:\ \text{finite background Rayleigh field}.
  \]
  Jump operator (outside minus inside), with $\eta:=(\tau;-\omega)$:
  \[
    \mathscr J_k[\eta;u_{\mathrm b}]
    =\begin{bmatrix}
      I+D^{(\kappa_{\mathrm e})}_\QP-D^{(\kappa_{\mathrm c})} & -S^{(\kappa_{\mathrm e})}_\QP+S^{(\kappa_{\mathrm c})}\\
      T^{(\kappa_{\mathrm e})}_\QP-T^{(\kappa_{\mathrm c})} & I-D'^{(\kappa_{\mathrm e})}_\QP+D'^{(\kappa_{\mathrm c})}
    \end{bmatrix}\eta
    +\begin{bmatrix}u_{\mathrm b}\\ \partial_\nu u_{\mathrm b}\end{bmatrix}_{\partial\Omega}.
  \]
  \begin{itemize}
    \item Transmission $\Leftrightarrow\mathscr J_k=0$: a second-kind system
      \cite{DominguezLyonTurc2016}. All kernel differences are at most log-singular,
      and the hypersingular part cancels in $T^{(\kappa_{\mathrm e})}-T^{(\kappa_{\mathrm c})}$.
    \item Pseudo-quasi-resonances \cite{HiptmairMoiolaSpence2022} are artifacts of the
      equation, not physical, and must be screened out.
  \end{itemize}
\end{frame}

\begin{frame}{Nystr\"om discretization and far-field extractors}
  \small
  \textbf{Nystr\"om.} Parametrize $\partial\Omega$ on $t_r=r/N_\eta$ and use Kress
  log-splitting \cite{Kress1991,HaoBarnettMartinssonYoung2014}; $n:=2N_\eta$ density
  unknowns, mesh size $h$.

  \medskip
  \textbf{Rayleigh coefficients of layer potentials.} If $\partial\Omega$ lies to the
  left of $x=x_p$, the layer potential is purely right-going for $x\ge x_p$, with
  amplitudes
  \[
    \phi_m^{\rightarrow}=\int_{\partial\Omega}\bigl(g_m^{\rightarrow}\omega
      +\partial_{\nu'}g_m^{\rightarrow}\tau\bigr)\dd s',\qquad
    g_m^{\rightarrow}(\mathbf z'):=\frac{\ii}{2\sqrt d\,\gamma_m}
      \ee^{-\ii\beta_my'+\ii\gamma_m(x_p-x')}.
  \]
  The left-going case (sources to the right) uses
  $g_m^{\leftarrow}:=\frac{\ii}{2\sqrt d\,\gamma_m}\ee^{-\ii\beta_my'+\ii\gamma_m(x'-x_p)}$.

  \medskip
  Discretizing with nodes $\mathbf z_j$, normals $\nu_j$, and weights $w_j$:
  \[
    \mathsf F^{\rightarrow/\leftarrow}_{x_p}
    :=\Bigl[\bigl(w_j\partial_{\nu_j}g_m^{\rightarrow/\leftarrow}(\mathbf z_j)\bigr)_{m,j}\ \
      \bigl(-w_jg_m^{\rightarrow/\leftarrow}(\mathbf z_j)\bigr)_{m,j}\Bigr]\in\CC^{K\times n},
    \qquad\text{acting on }\eta.
  \]
  The extractors are assembled directly from the spectral representation, not by
  port sampling followed by an FFT.
\end{frame}

\begin{frame}{One-cell scattering matrix}
  \small
  \begin{columns}[T]
    \begin{column}{0.42\textwidth}
      \centering
      \begin{tikzpicture}[x=0.9cm,y=0.9cm, every node/.style={font=\scriptsize},
          arr/.style={-{Latex[length=1.8mm]}, line width=0.7pt}]
        \fill[gray!6] (0,-1.2) rectangle (3,1.2);
        \draw[thick] (0,1.2) -- (3,1.2);
        \draw[thick] (0,-1.2) -- (3,-1.2);
        \draw[red!70!black, line width=1.1pt] (0,-1.2) -- (0,1.2);
        \draw[red!70!black, line width=1.1pt] (3,-1.2) -- (3,1.2);
        \draw[draw=blue!60!black, fill=blue!18] (1.5,0) circle (0.6);
        \node at (1.5,0) {$\Omega_+$};
        \node at (1.5,0.85) {$\eta$ on $\partial\Omega_+$};
        \draw[arr, green!45!black] (-1.3,0.5) -- (-0.1,0.5) node[midway, above] {$a_{\mathrm L}$ in};
        \draw[arr, orange!80!black] (-0.1,-0.5) -- (-1.3,-0.5) node[midway, below] {$b_{\mathrm L}$ out};
        \draw[arr, green!45!black] (4.3,-0.5) -- (3.1,-0.5) node[midway, below] {$b_{\mathrm R}$ in};
        \draw[arr, orange!80!black] (3.1,0.5) -- (4.3,0.5) node[midway, above] {$a_{\mathrm R}$ out};
        \node[below] at (0,-1.2) {$x_{\mathrm L}=x_+$};
        \node[below] at (3,-1.2) {$x_{\mathrm R}=x_++L$};
      \end{tikzpicture}
    \end{column}
    \begin{column}{0.56\textwidth}
      Incident field
      $u_{\mathrm b}=\sum_m\bigl(a_{\mathrm L,m}r_m^{x_{\mathrm L}}+b_{\mathrm R,m}l_m^{x_{\mathrm R}}\bigr)$.
      Solve $\mathscr J_k[\eta;u_{\mathrm b}]=0$ to get
      $\eta=\mathsf B_+(a_{\mathrm L};b_{\mathrm R})$.

      \medskip
      Layer potentials are left-going at $x_{\mathrm L}$ and right-going at
      $x_{\mathrm R}$:
      \[
        b_{\mathrm L}=\mathsf E(L)b_{\mathrm R}+\mathsf F^{\leftarrow}_{x_{\mathrm L}}\eta,\qquad
        a_{\mathrm R}=\mathsf E(L)a_{\mathrm L}+\mathsf F^{\rightarrow}_{x_{\mathrm R}}\eta.
      \]
    \end{column}
  \end{columns}

  \medskip
  \[
    \begin{bmatrix}b_{\mathrm L}\\ a_{\mathrm R}\end{bmatrix}
    =\begin{bmatrix}R_{\mathrm L}&T_{\mathrm{RL}}\\ T_{\mathrm{LR}}&R_{\mathrm R}\end{bmatrix}
    \begin{bmatrix}a_{\mathrm L}\\ b_{\mathrm R}\end{bmatrix},
    \qquad R_{\mathrm L},T_{\mathrm{RL}},T_{\mathrm{LR}},R_{\mathrm R}\in\CC^{K\times K}.
  \]
  The left cell $\mathcal C_-$ (ports $x_--L$, $x_-$) is treated the same way with its
  own $\Omega_-$. For BIE-based cell port maps see \cite{YuanLuAntoine2008,HuangLuLi2007}.
\end{frame}

\begin{frame}{Pencil, Floquet multipliers, and the decaying subspace}
  \small
  Port states $z_{\mathrm L}=(a_{\mathrm L};b_{\mathrm L})$ and
  $z_{\mathrm R}=(a_{\mathrm R};b_{\mathrm R})$:
  \[
    \mathsf E_{\mathrm L}z_{\mathrm L}=\mathsf E_{\mathrm R}z_{\mathrm R},\qquad
    \mathsf E_{\mathrm L}:=\begin{bmatrix}-R_{\mathrm L}&I\\ T_{\mathrm{LR}}&0\end{bmatrix},\quad
    \mathsf E_{\mathrm R}:=\begin{bmatrix}0&T_{\mathrm{RL}}\\ I&-R_{\mathrm R}\end{bmatrix}.
  \]
  Floquet solutions $z_{\mathrm R}=\lambda z_{\mathrm L}$ satisfy
  $\mathsf E_{\mathrm L}z=\lambda\mathsf E_{\mathrm R}z$, with $\lambda=\ee^{\ii\theta L}$
  a multiplier (not $k^2$). In the gap: no $|\lambda|=1$, and multipliers come in pairs
  $(\lambda,1/\lambda)$ \cite[Thm.~6.1]{JolyLiFliss2006}. We expect $K$ multipliers
  inside the unit circle and $K$ outside.

  \begin{block}{Ordered QZ (right side)}
    $\mathsf Q\mathsf E_{\mathrm L}\mathsf Z=\mathsf U_{\mathrm L}$ and
    $\mathsf Q\mathsf E_{\mathrm R}\mathsf Z=\mathsf U_{\mathrm R}$, with the $|\lambda|<1$
    cluster first. With $\mathsf Z_+:=$ the first $K$ columns,
    \[
      \mathsf P_+:=\mathsf U_{\mathrm R,11}^{-1}\mathsf U_{\mathrm L,11},\qquad
      \mathsf E_{\mathrm L}\mathsf Z_+=\mathsf E_{\mathrm R}\mathsf Z_+\mathsf P_+,\qquad
      \rho(\mathsf P_+)<1.
    \]
    $z_j:=\mathsf Z_+\mathsf P_+^jc$ gives consistent, geometrically decaying port states:
    $\mathsf E_{\mathrm L}z_j=\mathsf E_{\mathrm R}z_{j+1}$.
  \end{block}
\end{frame}

\begin{frame}{Left side, discrete Cauchy data space, and remarks}
  \footnotesize
  Left side: QZ of $(\mathsf E_{\mathrm R},\mathsf E_{\mathrm L})$ for $\mathcal C_-$ with the
  $|1/\lambda|<1$ cluster, giving
  $\mathsf E_{\mathrm R}\mathsf Z_-=\mathsf E_{\mathrm L}\mathsf Z_-\mathsf P_-$; then
  $\mathsf Z_-\mathsf P_-^jc$ are the right-port states of the cells going left.
  \begin{block}{Discrete Cauchy data space}
    \[
      \mathsf V_\pm:=\mathsf C_\pm\mathsf Z_\pm=\begin{bmatrix}\mathsf V_{\mathrm D,\pm}\\ \mathsf V_{\mathrm N,\pm}\end{bmatrix}\in\CC^{2K\times K},
      \qquad\mathcal G^h_\pm(k):=\operatorname{range}\mathsf V_\pm\approx\mathcal G_\pm(k)\
      \text{(Fourier coefficients, $|m|\le M$)}.
    \]
  \end{block}
  \begin{itemize}
    \item \textbf{Whole Schur subspace, not single eigenvectors.} This keeps multiple
      roots and non-normal coupling.
    \item \textbf{Basis invariance.} $\mathsf Z_\pm\mapsto\mathsf Z_\pm G$ gives
      $\mathsf P_\pm\mapsto G^{-1}\mathsf P_\pm G$ and $c\mapsto G^{-1}c$; the represented
      states are unchanged.
    \item \textbf{Failure detection.} The half-guide discretization is unusable at a
      parameter if the pencil is singular, if the selected cluster does not have
      dimension $K$, or if it is poorly separated from $|\lambda|=1$. Do not hide this
      by re-picking eigenvectors pointwise.
    \item \textbf{Non-uniqueness is not an issue.}
      $\mathsf E_{\mathrm L}z_{\mathrm L}=\mathsf E_{\mathrm R}z_{\mathrm R}$ does not
      determine $z_{\mathrm R}$ from $z_{\mathrm L}$ in general. The relation
      $\mathsf E_{\mathrm L}\mathsf Z_+=\mathsf E_{\mathrm R}\mathsf Z_+\mathsf P_+$ does
      not claim uniqueness; it \emph{selects} the decaying family, which is exactly
      what $\mathcal V_\pm(k)$ describes.
  \end{itemize}
\end{frame}

\begin{frame}{Center cell: representation theorem}
  \footnotesize
  \begin{theorem}[Center-cell representation]
    Assume no Wood anomaly, $\Omega_0=\varnothing$ or a $C^2$ simply connected domain
    compactly inside $\mathcal C_0$, and unique solvability of the full-plane outgoing
    transmission problem with swapped wavenumbers. If $u\in\mathcal U_0(k)$ is piecewise
    $H^2$ up to the boundary, then there exist $\tau,\omega$ and \emph{unique}
    $(a^-_{0,m}),(b^+_{0,m})$ such that
    \[
      u=\text{layer potentials on }\partial\Omega_0+u_{\mathrm{hom}},\qquad
      u_{\mathrm{hom}}=\sum_{m\in\ZZ}\bigl(a^-_{0,m}r_m^{x_-}+b^+_{0,m}l_m^{x_+}\bigr).
    \]
    The densities are also unique if the homogeneous complementary system has only
    the trivial solution.
  \end{theorem}
  \begin{itemize}
    \item $u_{\mathrm{hom}}$ is the part entering through the ports; the layer
      potentials are outgoing at both ports. $u_{\mathrm{hom}}\neq0$ is possible even
      when $\Omega_0=\varnothing$, so QP layer potentials alone are \emph{not}
      sufficient.
    \item \emph{Proof idea (appendix):} the Rayleigh coefficients come from the
      background bands next to the ports. An auxiliary problem with swapped
      wavenumbers defines $\tau,\omega$. Interior and exterior Green representations
      follow; port terms cancel because $w^+$ and $G_\QP$ are outgoing in the same
      direction, so their modal Wronskian vanishes.
  \end{itemize}
\end{frame}

\begin{frame}{Center cell: discrete unknowns and port data}
  \small
  Unknown $\xi:=(\eta;a^-_0;b^+_0)\in\CC^{n+2K}$ ($n=0$ if $\Omega_0=\varnothing$).
  The transmission equation $\mathscr J_k[\eta;u_{\mathrm{hom}}]=0$ reads
  \[
    J(k)\xi=0,\qquad J(k):=\bigl[J_\eta(k)\ \ \mathsf H_-(k)\ \ \mathsf H_+(k)\bigr]\in\CC^{n\times(n+2K)},
  \]
  \[
    (\mathsf H_-)_{jm}=\begin{bmatrix}1\\ \ii(\gamma_m\nu_{x,j}+\beta_m\nu_{y,j})\end{bmatrix}r_m^{x_-}(\mathbf z_j),\qquad
    (\mathsf H_+)_{jm}=\begin{bmatrix}1\\ \ii(-\gamma_m\nu_{x,j}+\beta_m\nu_{y,j})\end{bmatrix}l_m^{x_+}(\mathbf z_j).
  \]
  Amplitudes of the center field anchored at $\Gamma_\pm$, and the corresponding
  Cauchy-data coefficients:
  \[
    \mathsf Y_-\xi:=\begin{bmatrix}a^-_0\\ \mathsf F^{\leftarrow}_{x_-}\eta+\mathsf E(L_0)b^+_0\end{bmatrix},\quad
    \mathsf Y_+\xi:=\begin{bmatrix}\mathsf F^{\rightarrow}_{x_+}\eta+\mathsf E(L_0)a^-_0\\ b^+_0\end{bmatrix},\quad
    \Pi_\pm:=\mathsf C_\pm\mathsf Y_\pm=\begin{bmatrix}\Pi_{\mathrm D,\pm}\\ \Pi_{\mathrm N,\pm}\end{bmatrix}.
  \]
  \begin{alertblock}{}
    $J(k)\xi=0$ only imposes zero jumps \emph{at the Nystr\"om nodes}. The continuous
    jumps of the field built from finitely many coefficients and the exact kernels are
    re-evaluated in the error analysis.
  \end{alertblock}
\end{frame}

\begin{frame}{Discrete nonlinear eigenproblem}
  \small
  The discrete form of the center reduction: find $c_\pm\in\CC^K$ such that
  $J(k)\xi=0$ and $\Pi_\pm\xi=\mathsf V_\pm c_\pm$.
  \begin{block}{Cauchy-matching NEP}
    \[
      \mathsf M(k)\begin{bmatrix}\xi\\ c_-\\ c_+\end{bmatrix}
      :=\begin{bmatrix}J(k)\xi\\ \Pi_-\xi-\mathsf V_-c_-\\ \Pi_+\xi-\mathsf V_+c_+\end{bmatrix},
      \qquad\mathsf M(k)\in\CC^{(n+4K)\times(n+4K)},
      \qquad\text{find $k$ with $\mathsf M(k)$ singular.}
    \]
  \end{block}
  Since $\mathsf C_\pm$ is invertible, this is equivalent to amplitude matching
  $\mathsf Y_\pm\xi=\mathsf Z_\pm c_\pm$; the two forms differ by a row scaling.

  \medskip
  \textbf{DtN-reduced form (optional).} If $\mathsf V_{\mathrm D,\pm}$ is safely
  invertible in weighted coordinates, set
  $\widetilde\Lambda_\pm:=\mathsf V_{\mathrm N,\pm}\mathsf V_{\mathrm D,\pm}^{-1}$ and eliminate $c_\pm$:
  \[
    \mathsf M_{\mathrm{DtN}}(k)\xi:=
    \bigl(J(k)\xi;\ \Pi_{\mathrm N,-}\xi-\widetilde\Lambda_-\Pi_{\mathrm D,-}\xi;\
      \Pi_{\mathrm N,+}\xi-\widetilde\Lambda_+\Pi_{\mathrm D,+}\xi\bigr),\qquad\text{order }n+2K.
  \]
  A moderate $\operatorname{cond}(\mathsf V_{\mathrm D,\pm})$ alone does not guarantee
  safety. Failure of the Dirichlet coordinates is the discrete analogue of the
  Dirichlet exceptional values. \textbf{Default: $\mathsf M(k)$.}
\end{frame}

\begin{frame}{Candidate localization and output}
  \footnotesize
  With invertible row/column weights $W_r(k),W_c(k)$ (balancing interface rows,
  value/normal-derivative coefficients, and unknown scales):
  \[
    s(k):=\frac{\sigma_{\min}\bigl(W_r\mathsf M(k)W_c\bigr)}{\sigma_{\max}\bigl(W_r\mathsf M(k)W_c\bigr)}.
  \]
  \begin{itemize}
    \item Coarse scan on $(\sqrt{g_-},\sqrt{g_+})$, then rule-based refinement around
      isolated dips, giving $\widehat k$.
    \item Checks: QZ cluster separation and dimension, unweighted residual, and
      consistency across neighboring discretization levels $(h,M)$. These screen out
      pseudo-quasi-resonances and other artifacts.
  \end{itemize}

  \begin{block}{Candidate output}
    $\widehat y$ is the unit right singular vector for $\sigma_{\min}$ of
    $W_r\mathsf M(\widehat k)W_c$, and
    \[
      (\xi;c_-;c_+):=W_c\widehat y\qquad(\text{defined up to a nonzero scalar}).
    \]
    Inputs to the error analysis: $\widehat k$, $(\xi,c_\pm)$, $(h,M)$,
    $\mathsf P_\pm$, and $\mathsf Z_\pm$. A small $s(\widehat k)$ only says the discrete
    system is nearly solvable; it gives \textbf{no} distance from $\widehat k^2$ to
    $\sigma(A(\beta))$.
  \end{block}
\end{frame}

% ======================================================================
\section{Spectral Distance Estimate}

\begin{frame}{Strategy}
  \small
  Let $\mu:=\widehat k^2$.
  \begin{enumerate}
    \item Build a \textbf{raw field} $\widehat u$ on the whole strip from the candidate
      coefficients and the \emph{exact} kernels. It solves Helmholtz exactly in each
      homogeneous piece and is exactly $\beta$-quasi-periodic. It leaves small jumps on
      material and artificial interfaces, so $\widehat u\notin D(A(\beta))$.
    \item Collect all jumps into the interface defect $\eps$, and set
      $U:=\lVert\widehat u\rVert_{\mathcal H}$.
    \item A trace estimate for the positive-shift resolvent and a Green identity give
      \[
        \dist(\mu,\sigma(A(\beta)))\le\frac{2\mu Q}{U-Q},\qquad Q=C\eps,
      \]
      where $C$ depends only on geometry, material bounds, and $\mu$.
  \end{enumerate}

  \begin{block}{What the theorem does not need}
    No exact discrete null vector, no uniqueness of $\widehat u$, and no resolvent
    solve. The details of the candidate algorithm enter only through $\eps$ and $U$.
  \end{block}
\end{frame}

\begin{frame}{Raw field: construction}
  \footnotesize
  \textbf{Definition vs.\ evaluation parameters.} $(h\sim N_\eta, M, \text{candidate vector})$
  fix $\widehat u$. The parameters $(N_q, M_{\mathrm{ev}}\ge M, N_{\mathrm{cell}},
  \text{kernel accuracy})$ only control how accurately $\eps,U$ are computed.
  Both quantities are 1-homogeneous in the scalar multiple, so the bound is
  normalization-free.

  \medskip
  \textbf{Continuous density.} Trigonometric interpolation of the $N_\eta$ samples:
  \[
    \eta_m:=\frac1{N_\eta}\sum_{r=0}^{N_\eta-1}\eta_r\ee^{-2\pi\ii mr/N_\eta},\qquad
    \eta(t):=\sum_{m=-N_\eta/2}^{N_\eta/2-1}\eta_m\ee^{2\pi\ii mt}.
  \]
  Nystr\"om sums $\sum_rw_rG(\cdot,\mathbf z_r)\omega_r$ cannot serve as the field: they
  carry point-source and dipole singularities and are neither piecewise $H^2$ nor $L^2$.

  \medskip
  \textbf{Pieces.}
  \begin{itemize}
    \item Center: layer potentials with the continuous density, plus $u_{\mathrm{hom}}$
      ($|m|\le M$) at $k=\widehat k$.
    \item Lead cell: $u^{\mathrm{cell}}_\pm[a_{\mathrm L},b_{\mathrm R}]$ from
      $\mathsf B_\pm$ and the continuous density. With $\mathsf Z_\pm v=(a_\pm(v);b_\pm(v))$:
      \[
        u_+[v]:=u_+^{\mathrm{cell}}\bigl[a_+(v),b_+(\mathsf P_+v)\bigr],\quad
        u_-[v]:=u_-^{\mathrm{cell}}\bigl[a_-(\mathsf P_-v),b_-(v)\bigr].
      \]
    \item Half-guides: $\widehat u(\mathbf z\pm jLe_x):=u_\pm[\mathsf P_\pm^jc_\pm](\mathbf z)$
      for $\mathbf z\in\mathcal C_\pm$, $j\ge0$.
  \end{itemize}
\end{frame}

\begin{frame}{Raw field: consistency and properties}
  \footnotesize
  \textbf{Self-consistency on $|m|\le M$.} With $z_j=\mathsf Z_+\mathsf P_+^jc_+=(a_j;b_j)$,
  the relation $\mathsf E_{\mathrm L}z_j=\mathsf E_{\mathrm R}z_{j+1}$ gives
  $b_j=R_{\mathrm L}a_j+T_{\mathrm{RL}}b_{j+1}$ and
  $a_{j+1}=T_{\mathrm{LR}}a_j+R_{\mathrm R}b_{j+1}$. So cell $j$ with incident
  $(a_j,b_{j+1})$ emits exactly $b_j$ and $a_{j+1}$, and both sides of each shared port
  carry the amplitudes $(a_{j+1},b_{j+1})$.

  The remaining mismatch comes from modes with $|m|>M$, from the gap between the
  continuous and the discrete densities, and from $\sigma_{\min}\neq0$. The field
  decays geometrically cell by cell since $\rho(\mathsf P_\pm)<1$.

  \begin{block}{Properties used by the theorem}
    \begin{itemize}
      \item[(P1)] $(-\Delta-\mu\rho)\widehat u=0$ in each homogeneous piece (exact kernels
        and Rayleigh waves);
      \item[(P2)] value and $\partial_y$ traces are exactly $\beta$-quasi-periodic;
      \item[(P3)] locally $H^2$ up to the boundary of each piece (assumed);
      \item[(P4)] $\widehat u\in\mathcal H$ (from propagation decay; checked via the mass sum).
    \end{itemize}
    But $\widehat u\notin H^1_\beta(B)$: the value jumps are nonzero in general.
  \end{block}
\end{frame}

\begin{frame}{Interface jumps and the two input quantities}
  \footnotesize
  Interface family (each shared curve counted once; the top/bottom boundaries cancel
  by quasi-periodicity):
  \[
    \mathscr E:=\mathscr E_{\mathrm{mat}}\dot\cup\mathscr E_{\mathrm{art}},\qquad
    \mathscr E_{\mathrm{mat}}:=\{\partial\Omega_0\}\cup\{\partial\Omega_{j,\pm}\}_{j\ge0},\qquad
    \mathscr E_{\mathrm{art}}:=\{\Gamma_\pm\}\cup\{\Gamma_{j,\pm}\}_{j\ge0}.
  \]
  Orientation: $\nu_\Gamma$ points from side 1 to side 2. Side 1 is the inclusion
  interior on material interfaces and the cell nearer the center on artificial
  interfaces.
  \[
    \jump{\widehat u}_\Gamma:=\widehat u^{(2)}-\widehat u^{(1)},\qquad
    \jump{\partial_\nu\widehat u}_\Gamma:=\partial_{\nu_\Gamma}\widehat u^{(1)}-\partial_{\nu_\Gamma}\widehat u^{(2)}
  \]
  (signs chosen so that both terms enter the Green identity with $+$).
  \begin{itemize}
    \item On material interfaces, $(\jump{\widehat u},-\jump{\partial_\nu\widehat u})$ is
      exactly $\mathscr J_k$ applied with \emph{continuous} operators and densities.
      $J(\widehat k)\xi\approx0$ only controls it at the nodes.
    \item The port rows of $\mathsf M(\widehat k)$ are approximately the $|m|\le M$ jump
      coefficients on $\Gamma_\pm$.
  \end{itemize}
  \begin{block}{Inputs ($\ell_*>0$ an arbitrary reference length)}
    \[
      \eps^2:=\sum_{\Gamma\in\mathscr E}\Bigl(\lVert\jump{\partial_\nu\widehat u}_\Gamma\rVert_{L^2(\Gamma)}^2
        +\ell_*^{-2}\lVert\jump{\widehat u}_\Gamma\rVert_{L^2(\Gamma)}^2\Bigr),\qquad
      U:=\lVert\widehat u\rVert_{\mathcal H}.
    \]
  \end{block}
\end{frame}

\begin{frame}{Trace estimate for the positive-shift resolvent}
  \footnotesize
  Geometry: normal neighborhoods $\boldsymbol r_\Gamma(s)+t\nu_\Gamma(s)$, $|t|<\delta_{\mathrm{col}}$,
  injective; curvature bound $\kappa_*$; overlap multiplicity $N_{\mathrm{ov}}$;
  \[
    \vartheta_{\mathrm{col}}:=1-\delta_{\mathrm{col}}\kappa_*>0,\qquad
    C_0:=\frac{N_{\mathrm{ov}}}{\vartheta_{\mathrm{col}}\delta_{\mathrm{col}}},\qquad
    C_1:=\frac{N_{\mathrm{ov}}\delta_{\mathrm{col}}}{2\vartheta_{\mathrm{col}}}.
  \]
  \begin{proposition}[All-interface trace bound]
    For $v\in D(A(\beta))$ and $\mu>0$,
    \begin{align*}
      T_{\ell_*}(v)^2&:=\sum_{\Gamma\in\mathscr E}\Bigl(\lVert v\rVert_{L^2(\Gamma)}^2
        +\ell_*^2\lVert\partial_{\nu_\Gamma}v\rVert_{L^2(\Gamma)}^2\Bigr)
      \le C_{\mathrm{tr}}^2\lVert(A(\beta)+\mu)v\rVert_{\mathcal H}^2,\\
      C_{\mathrm{tr}}^2&:=\frac{C_0}{\rin\mu^2}+\frac{C_1+\ell_*^2C_0}{4\mu}+\ell_*^2C_1\rax.
    \end{align*}
  \end{proposition}
  These conditions do not follow from simple connectedness alone. For certified
  bounds, neighborhood injectivity, separation of interfaces and their periodic
  copies, and the overlap multiplicity must be verified.

  \medskip
  Notation: $C:=C_{\mathrm{tr}}$ and $Q:=C\eps$. $C$ is \emph{not} fitted from residuals.
\end{frame}

\begin{frame}{Trace estimate: proof ingredients}
  \small
  Let $v\in D(A(\beta))$ and $f:=(A(\beta)+\mu)v$.
  \begin{itemize}
    \item Global $H^2$ via Fourier: $\lVert D^2v\rVert_{L^2(B)}=\lVert\Delta v\rVert_{L^2(B)}$.
    \item Spectral calculus: $\lVert v\rVert^2\le\lVert f\rVert^2/(\rin\mu^2)$,
      $\lVert\nabla v\rVert^2\le\lVert f\rVert^2/(4\mu)$,
      $\lVert D^2v\rVert^2\le\rax\lVert f\rVert^2$, where $f=(A+\mu)v$.
    \item 1D trace inequality
      $|g(0)|^2\le\delta^{-1}\!\int_{-\delta}^{\delta}|g|^2+\frac\delta2\int_{-\delta}^{\delta}|g'|^2$
      along normals, applied to $v$ and $\partial_\nu v$ ($|g'|\le|D^2v|$), then
      integrated using the area element $1-t\kappa\ge\vartheta_{\mathrm{col}}$ and summed
      with $N_{\mathrm{ov}}$:
      \begin{align*}
        \textstyle\sum_\Gamma\lVert v\rVert_{L^2(\Gamma)}^2
        &\le C_0\lVert v\rVert_{L^2(B)}^2+C_1\lVert\nabla v\rVert_{L^2(B)}^2,\\
        \textstyle\sum_\Gamma\lVert\partial_{\nu_\Gamma}v\rVert_{L^2(\Gamma)}^2
        &\le C_0\lVert\nabla v\rVert_{L^2(B)}^2+C_1\lVert D^2v\rVert_{L^2(B)}^2.
      \end{align*}
  \end{itemize}
  The trace bound is what controls $z=(A(\beta)-\mu)^{-1}\widehat u$ on all interfaces
  in the main proof.
\end{frame}

\begin{frame}{Main theorem}
  \begin{theorem}[Spectral distance from raw interface defects]
    Assume the geometric conditions above, $\mu>0$, a nonzero $\widehat u$ satisfying
    (P1)--(P4), and $\eps<\infty$. If $Q<U$, then
    \[
      \boxed{\ \dist\bigl(\mu,\sigma(A(\beta))\bigr)\le\frac{2\mu Q}{U-Q},\qquad Q=C_{\mathrm{tr}}\,\eps.\ }
    \]
    This is a distance to the \emph{whole} spectrum; no knowledge of
    $\sigma_{\mathrm{ess}}$ or the gap is needed.
  \end{theorem}
  \small
  \begin{itemize}
    \item The bound is invariant under scaling of $\widehat u$ and tends to zero linearly,
      $\sim2\mu C\eps/U$.
    \item It holds for every $\ell_*>0$, which can balance the two jump types.
    \item \emph{Volume defect variant:} if (P1) fails with
      $r=(-\Delta-\mu\rho)\widehat u$ and
      $\varepsilon_{\mathrm{vol}}=\lVert\rho^{-1}r\rVert_{\mathcal H}$, the bound becomes
      $(2\mu Q+\varepsilon_{\mathrm{vol}})/(U-Q)$. For the raw field,
      $\varepsilon_{\mathrm{vol}}=0$.
  \end{itemize}
\end{frame}

\begin{frame}{Proof sketch}
  \footnotesize
  Let $\delta:=\dist(\mu,\sigma(A(\beta)))>0$ and $z:=(A(\beta)-\mu)^{-1}\widehat u$
  (used only in the proof, never computed). Write $u=\widehat u$.
  \begin{enumerate}
    \item \textbf{Piecewise identity.} $\Delta u=-\mu\rho u$ and
      $-\Delta z-\mu\rho z=\rho u$ give $\rho|u|^2=\overline z\Delta u-u\Delta\overline z$.
      Green's second identity applies on each piece.
    \item \textbf{Merge shared interfaces.} $z\in H^2(B)$ has equal traces on both sides,
      so
      $(\overline z\partial_\nu u^{(1)}-u^{(1)}\partial_\nu\overline z)
      -(\overline z\partial_\nu u^{(2)}-u^{(2)}\partial_\nu\overline z)
      =\jump{\partial_\nu u}\overline z+\jump{u}\partial_\nu\overline z$.
      The top and bottom boundaries cancel.
    \item \textbf{Cutoff $\chi_R$.} The extra terms $I_R\to0$ (since $u\in L^2$, $z\in H^2$)
      and $|J_R|\le\lVert\chi_R'\rVert_\infty\ell_*\eps T_{\ell_*}(z)\to0$. The latter uses
      the trace bound, which is needed exactly to remove the cutoff.
    \item \textbf{Identity.}
      $\displaystyle U^2=\sum_{\Gamma\in\mathscr E}\int_\Gamma\bigl(\jump{\partial_\nu u}_\Gamma\overline z
        +\jump{u}_\Gamma\partial_{\nu_\Gamma}\overline z\bigr)\dd s\le\eps\,T_{\ell_*}(z).$
    \item \textbf{Resolvent bound.} $\lVert z\rVert\le U/\delta$ and
      $(A+\mu)z=u+2\mu z$, so $T_{\ell_*}(z)\le C(1+2\mu/\delta)U$. Hence
      \[
        U^2\le Q\Bigl(1+\frac{2\mu}{\delta}\Bigr)U
        \quad\Longrightarrow\quad\delta(U-Q)\le2\mu Q.
      \]
  \end{enumerate}
\end{frame}

\begin{frame}{Computing $\eps$: material interfaces}
  \small
  \textbf{Material interface} $\Gamma=\partial\Omega$, targets $t_q=q/N_q$:
  \begin{itemize}
    \item local terms $\tau(t_q),\omega(t_q)$ come directly from the trigonometric
      interpolant;
    \item integral terms: zero-pad the density to $N_q\ge N_\eta$ points (exact, since
      the degree is $<N_\eta/2$), then apply Kress splitting. The kernel differences
      are at most log-singular, so no hypersingular quadrature is needed;
    \item outer periodic trapezoid rule for the $L^2$ norm:
      \[
        \int_\Gamma\bigl(|\jump{\partial_\nu\widehat u}|^2+\ell_*^{-2}|\jump{\widehat u}|^2\bigr)\dd s
        \approx\frac1{N_q}\sum_{q=0}^{N_q-1}\bigl(|\jump{\partial_\nu\widehat u}(\mathbf z_q)|^2
          +\ell_*^{-2}|\jump{\widehat u}(\mathbf z_q)|^2\bigr)|\boldsymbol r'(t_q)|.
      \]
  \end{itemize}
  Two quadrature levels: the inner one evaluates layer potentials at targets, the outer
  one computes the $L^2$ norm. $N_q\gg N_\eta$ is allowed without changing the density.
  The background field $u_{\mathrm b}$ and its normal derivative are evaluated directly.
\end{frame}

\begin{frame}{Computing $\eps$: artificial interfaces}
  \small
  \textbf{Artificial interface} $\Gamma=\{x_\Gamma\}\times I_d$, using Parseval in $\psi_m$:
  \[
    \int_\Gamma\bigl(|\jump{\partial_\nu\widehat u}|^2+\ell_*^{-2}|\jump{\widehat u}|^2\bigr)
    =\sum_{m\in\ZZ}\bigl(|\jump{\partial_\nu\widehat u}_m|^2+\ell_*^{-2}|\jump{\widehat u}_m|^2\bigr).
  \]
  On each side the field near $\Gamma$ is Rayleigh waves plus QP layer potentials whose
  sources lie on one side. Rayleigh waves give amplitudes directly; layer potentials
  give $\phi_m$ (value) and $\pm\ii\gamma_m\phi_m$ ($\partial_x$).

  \medskip
  Compute $|m|\le M_{\mathrm{ev}}$ explicitly; the rest cannot be set to zero. Take
  $M_{\mathrm{ev}}$ so that $|m|>M_{\mathrm{ev}}$ is evanescent (layer-potential part
  only), and let $\delta_\Gamma$ be the horizontal distance to the source curves:
  \begin{align*}
    |\phi_m|\le\Phi_m&:=\frac{\ee^{-|\gamma_m|\delta_\Gamma}}{2\sqrt d|\gamma_m|}
    \Bigl(\lVert\omega\rVert_{L^1}+(|\gamma_m|^2+\beta_m^2)^{1/2}\lVert\tau\rVert_{L^1}\Bigr),\\
    \sum_{|m|>M_{\mathrm{ev}}}\bigl(|\jump{\partial_\nu\widehat u}_m|^2+\ell_*^{-2}|\jump{\widehat u}_m|^2\bigr)
    &\le2\sum_{|m|>M_{\mathrm{ev}}}(|\gamma_m|^2+\ell_*^{-2})\bigl((\Phi_m^{(1)})^2+(\Phi_m^{(2)})^2\bigr).
  \end{align*}
  The tail decays like $\exp(-4\pi|m|\delta_\Gamma/d)$ times a polynomial, giving a
  closed-form geometric bound.
\end{frame}

\begin{frame}{Computing $\eps$: infinitely many cells}
  \small
  By translation structure, the jumps on $\partial\Omega_{j,\pm}$ and $\Gamma_{j,\pm}$
  depend only on the state $\mathsf P_\pm^jc_\pm$:
  \[
    \eps^2=b_0+\sum_\pm\sum_{j=0}^\infty b_\pm(\mathsf P_\pm^jc_\pm),\qquad
    b_\pm(v)=\lVert\mathcal L_\pm v\rVert^2\ \ (\mathcal L_\pm\text{ linear}).
  \]
  Here $b_0$ collects $\partial\Omega_0$ and $\Gamma_\pm$, and $b_\pm(v)$ collects
  $\partial\Omega_\pm$ and the shared port $\Gamma_{0,\pm}$.

  \begin{block}{Propagation tail}
    $b_\pm(v)\le\Theta_\pm\lVert v\rVert_2^2$ with $\Theta_\pm:=\sum_{s=1}^Kb_\pm(\boldsymbol e_s)=\lVert\mathcal L_\pm\rVert_{\mathrm{HS}}^2$. If
    $\lVert\mathsf P_\pm^p\rVert_2\le r_\pm<1$ for some $p$:
    \[
      \sum_{j\ge N_{\mathrm{cell}}}b_\pm(\mathsf P_\pm^jc_\pm)\le
      \frac{\Theta_\pm\sum_{l=0}^{p-1}\lVert\mathsf P_\pm^l\rVert_2^2}{1-r_\pm^2}
      \bigl\lVert\mathsf P_\pm^{N_{\mathrm{cell}}}c_\pm\bigr\rVert_2^2.
    \]
    $\mathsf P_\pm$ is generally non-normal, so $\lVert\mathsf P_\pm\rVert_2>1$ is possible
    even when $\rho(\mathsf P_\pm)<1$; then $p>1$ is needed. The tail cannot be dropped
    just because the last term is small.
  \end{block}
\end{frame}

\begin{frame}{Field norm and procedure}
  \small
  \textbf{Mass.} $U^2=\int_{\mathcal C_0}\rho|\widehat u|^2+\sum_\pm\sum_jm_\pm(\mathsf P_\pm^jc_\pm)$,
  and the same tail argument verifies (P4). The theorem only needs a lower bound
  $\bigl(\int_D\rho|\widehat u|^2\bigr)^{1/2}\le U$ for a subregion $D$ away from
  interfaces, where plain trapezoid sums suffice. $D$ must carry most of the mass, or
  $Q<U$ fails.

  \medskip
  \textbf{Procedure}
      \begin{enumerate}
        \item get $(\xi,c_-,c_+)$ and fix the multiple;
        \item compute $b_0$;
        \item sum $b_\pm(\mathsf P_\pm^jc_\pm)$ for $j<N_{\mathrm{cell}}$, then add the
          tail using $\Theta_\pm$ and $r_\pm$; this gives $\eps$;
        \item compute a lower bound for $U$ on $D$;
        \item compute $C$ and $Q=C\eps$, check $Q<U$, and apply the theorem.
      \end{enumerate}
      Increase $N_q$, $M_{\mathrm{ev}}$, and $N_{\mathrm{cell}}$ until all three
      contributions (interface integrals, modal tails, propagation tails) are stable.
      This is a numerical \emph{estimate}: refinement stability is not a certificate.
\end{frame}

\begin{frame}{Certified bounds and interpretation}
  \small
      \textbf{One-sided inputs}
      \[
        C\le\overline C,\quad\eps\le\overline\varepsilon_\partial,\quad
        0<\underline U\le U,\quad
        \underline D:=\underline U-\overline C\,\overline\varepsilon_\partial>0.
      \]
      All three bounds use the same candidate, the same multiple, and the same
      densities. They must account for:
      \begin{itemize}
        \item geometric checks and outward rounding for $\overline C$;
        \item kernel, inner/outer quadrature, and rounding errors, plus the modal and
          propagation tails for $\overline\varepsilon_\partial$;
        \item quadrature and rounding errors for $\underline U$.
      \end{itemize}
      \[
        R_{\mathrm{spec}}:=\frac{2\mu\overline C\,\overline\varepsilon_\partial}{\underline D},
        \qquad
        I_{\mathrm{spec}}:=[\mu-R_{\mathrm{spec}},\mu+R_{\mathrm{spec}}],
        \qquad
        I_{\mathrm{spec}}\cap\sigma(A(\beta))\neq\varnothing.
      \]

  \textbf{Interpretation.} $I_{\mathrm{spec}}$ meets the spectrum, but gives no index,
  multiplicity, or uniqueness. If $I_{\mathrm{spec}}$ is \emph{separately} shown to lie
  in an open interval disjoint from $\sigma_{\mathrm{ess}}$, it contains at least one
  isolated, finite-multiplicity eigenvalue: a guided mode.
\end{frame}

% ======================================================================
\section{Summary}

\begin{frame}{Summary}
  \small
  \centering
  \begin{tikzpicture}[node distance=0.42cm,
      box/.style={draw=black!60, rounded corners=2pt, fill=blue!5, align=center,
        font=\scriptsize, minimum height=0.95cm, text width=3.2cm},
      arr/.style={-{Latex[length=1.8mm]}, line width=0.6pt}]
    \node[box] (a) {model $A(\beta)$\\ $\sigma_{\mathrm{ess}}$, gap $G_\beta$};
    \node[box, right=of a] (b) {center reduction\\ $\mathscr T_\pm u_0\in\mathcal G_\pm(k)$};
    \node[box, right=of b] (c) {BIE + Rayleigh ports\\ $\mathsf P_\pm$, $\mathsf V_\pm$ via QZ};
    \node[box, below=0.7cm of c] (d) {NEP $\mathsf M(k)$\\ candidate $\widehat k$, $(\xi,c_\pm)$};
    \node[box, left=of d] (e) {raw field $\widehat u$\\ jumps $\eps$, mass $U$};
    \node[box, left=of e] (f) {$\dist(\mu,\sigma)\le\dfrac{2\mu Q}{U-Q}$\\ certified $I_{\mathrm{spec}}$};
    \draw[arr] (a) -- (b);
    \draw[arr] (b) -- (c);
    \draw[arr] (c) -- (d);
    \draw[arr] (d) -- (e);
    \draw[arr] (e) -- (f);
  \end{tikzpicture}

  \vspace{0.6em}
  \raggedright
  \begin{itemize}
    \item The Cauchy-data reduction is exact for every $k>0$; DtN is the graph special case.
    \item The discretization mirrors the continuous reduction: $\mathcal G_\pm\to
      \operatorname{range}\mathsf C_\pm\mathsf Z_\pm$.
    \item The candidate algorithm enters the error bound only through $\eps$ and $U$. The
      bound holds for the whole spectrum, and the gap is used only to interpret the
      spectral point as a guided mode.
  \end{itemize}

  \textbf{Next:} numerical results (convergence of $\eps$, tail bounds, certified
  intervals).
\end{frame}

\begin{frame}[allowframebreaks]{References}
  \scriptsize
  \bibliographystyle{plain}
  \bibliography{../../../ref/reference}
\end{frame}

\end{document}
